\subsection{Drell--Yan 过程：从颜色平均到强子截面}
\label{sec:qcd-drell-yan-applications}

DIS 用轻子探测一个强子，强子碰撞则同时含有两个未知的长距离初态。Drell--Yan 过程 $H_AH_B\to\ell^-\ell^++X$ 的末态轻子不参与强子化，因此适合检验从 DIS 提取的 PDF 能否描述另一类实验。其一般起点是对未分辨强子末态 $X$ 求和的电流矩阵元平方；在硬尺度远大于强子尺度时，再把这一包容量因子化。因子化的对象是截面，不能把两个强子当作确定的单个夸克态。

令 $P_A,P_B$ 为入射强子四动量，$K=k_-+k_+$ 为轻子对总四动量。沿用物理动量量纲，定义
\begin{equation}
 s=(P_A+P_B)^2,\qquad \mathcal Q^2=K^2>0,\qquad
 Y=\frac12\ln\frac{K^0+K^3}{K^0-K^3},\qquad
 \tau_{\rm DY}=\frac{\mathcal Q^2}{s}.
 \label{eq:dy-observables}
\end{equation}
$Y$ 是强子质心系中沿 $H_A$ 束流方向的轻子对快度；$c\mathcal Q$ 是轻子对不变质量对应的能量。这里的 $\mathcal Q^2$ 为类时正数，DIS 则定义为 $-q^2$，两者都表示正的硬尺度平方。

在远离强子阈值和极端小动量分数、并对强子辐射与轻子对横向动量作包容积分时，领先幂共线因子化写为
\begin{equation}
 \frac{\dd^2\sigma}{\dd\mathcal Q^2\dd Y}
 =\sum_{i,j}\int_0^1\dd x_1\int_0^1\dd x_2\,
 f_{i/H_A}(x_1,\mu_{\rm F})f_{j/H_B}(x_2,\mu_{\rm F})
 \frac{\dd^2\hat\sigma_{ij}}{\dd\mathcal Q^2\dd Y}
 +\Delta_{\rm power}.
 \label{eq:dy-factorization}
\end{equation}
短距离截面 $\hat\sigma_{ij}$ 可逐阶微扰计算；$\Delta_{\rm power}$ 的典型相对量级为 $\Lambda_{\rm had}^2/\mathcal Q^2$，其中 $\Lambda_{\rm had}$ 是强子动量尺度。在运动学端点，这一估计一般不均匀。高阶硬系数含额外辐射卷积，不能继续把 $x_1,x_2$ 当作由 $\mathcal Q,Y$ 唯一固定的数值。

先取纯虚光子交换、无质量外部夸克与轻子，并保留强耦合的最低阶。部分子运动学为 $\hat s=(p_1+p_2)^2$、$\hat t=(p_1-k_-)^2$、$\hat u=(p_1-k_+)^2$，满足 $\hat s+\hat t+\hat u=0$。无量纲电磁耦合 $g_{\rm em}\equiv\sqrt{4\pi\alpha}$ 与 SI 电荷 $e$ 的关系是 $\alpha=e^2/(4\pi\varepsilon_0\hbar c)$。对初态自旋和颜色平均后，
\begin{equation}
 \overline{|\mathcal M_{q_f\bar q_f}|^2}
 =\frac{2g_{\rm em}^4Q_f^2}{N_c}
 \frac{\hat t^2+\hat u^2}{\hat s^2},\qquad
 \hat\sigma_{q_f\bar q_f}
 =\frac{4\pi\alpha^2\hbar^2Q_f^2}{3N_c\hat s}.
 \label{eq:dy-born-partonic}
\end{equation}
$Q_f$ 是以 $e$ 为单位的夸克电荷。这里 $\hbar^2/\hat s$ 的单位是面积；若用能量平方 $c^2\hat s$，同一前因子写成 $(\hbar c)^2/(c^2\hat s)$。

在 Born 共线运动学中，$\hat s=x_1x_2s$，$Y=\frac12\ln(x_1/x_2)$，因此
\begin{equation}
 x_1=\sqrt{\tau_{\rm DY}}\ee^Y,\qquad
 x_2=\sqrt{\tau_{\rm DY}}\ee^{-Y},\qquad
 |Y|\leq\ln\frac{\sqrt s}{\mathcal Q}.
 \label{eq:dy-born-fractions}
\end{equation}
定义这一运动学点的无量纲味权重
\begin{equation*}
 \mathcal W_f^{AB}
 =f_{f/H_A}(x_1)f_{\bar f/H_B}(x_2)
 +f_{\bar f/H_A}(x_1)f_{f/H_B}(x_2),
\end{equation*}
其中同一 $\mu_{\rm F}$ 自变量暂略去。两项对应夸克来自两个不同束流的两种可能性。可测截面为
\begin{equation}
 \frac{\dd^2\sigma_{\rm Born}}{\dd\mathcal Q^2\dd Y}
 =\frac{4\pi\alpha^2\hbar^2}{3N_c s\mathcal Q^2}
 \sum_f Q_f^2\mathcal W_f^{AB}.
 \label{eq:dy-born-hadronic}
\end{equation}
对于相同非极化束流，交换 $A,B$ 等价于 $Y\to-Y$，因而分布必须关于 $Y=0$ 对称。纯光子近似还要求远离 $Z$ 极点；接近电弱尺度时，应将前文的手征 $\gamma$--$Z$ 振幅放入同一因子化公式。相关实验正是通过质量和快度的双微分分布检验 PDF 与高阶硬系数\cite{CMS2015DrellYan}。

\begin{exbox}[质量分箱、快度与事件数]
设强子质心能量为 $E_{\rm cm}=c\sqrt s=\SI{200}{\giga\electronvolt}$，轻子对能量质量为 $M_{\ell\ell}=c\mathcal Q=\SI{10}{\giga\electronvolt}$。在 $Y=1$ 时，
\begin{equation*}
 \tau_{\rm DY}=0.0025,\qquad x_1=0.1359,\qquad x_2=0.01839.
\end{equation*}
这说明同一质量分箱的正快度事件同时取样一个较大 $x$ 和一个较小 $x$。作单位演算时，假设该点给定的 PDF 味和为 $\sum_fQ_f^2\mathcal W_f^{AB}=20$，并取 $\alpha=1/137$。利用 $\dd\mathcal Q^2/\dd M_{\ell\ell}=2M_{\ell\ell}/c^2$，
\begin{align*}
 \frac{\dd^2\sigma}{\dd M_{\ell\ell}\dd Y}
 &=\frac{8\pi\alpha^2(\hbar c)^2}{9E_{\rm cm}^2M_{\ell\ell}}
 \sum_fQ_f^2\mathcal W_f^{AB}\\
 &\simeq\SI{2.90}{\pico\barn\per\giga\electronvolt}.
\end{align*}
这里使用 $(\hbar c)^2=0.38938\,\mathrm{mb\,GeV^2}$。若足够窄的分箱取 $\Delta M_{\ell\ell}=\SI{0.1}{\giga\electronvolt}$、$\Delta Y=0.1$，局部常值估计给 $\Delta\sigma\simeq\SI{0.0290}{\pico\barn}$。积分亮度为 $\mathcal L_{\rm int}=\SI{1000}{\per\pico\barn}$、平均接收效率为 $\mathcal Q^{\rm det}=0.50$ 时，期望事件数约为 $14.5$。

事件数是期望值，可以不是整数。这一计算只演示 Jacobian、单位和亮度换算；给定味和不是实际 PDF 拟合，高阶 QCD 修正、分箱内变化与探测响应也未计入。对宽分箱应积分微分截面与响应函数，不能再用中心点乘分箱宽度。
\end{exbox}

\begin{derivation}[颜色平均、自旋迹与两次 Jacobian]
虚光子不改变颜色。取夸克和反夸克颜色为 $i,j$，其振幅含 $\delta_{ij}$，故
\begin{equation*}
 \frac1{N_c^2}\sum_{i,j}|\delta_{ij}|^2=\frac1{N_c}.
\end{equation*}
它与 $e^+e^-\to q\bar q$ 的末态颜色求和 $N_c$ 不同；交叉振幅并不会自动把初态平均一起交叉。
定义
\begin{equation*}
 T^{\mu\nu}(a,b)=\Tr(\slashed a\gamma^\mu\slashed b\gamma^\nu)
 =4(a^\mu b^\nu+a^\nu b^\mu-\eta^{\mu\nu}a\cdot b).
\end{equation*}
两迹的指标缩并为
\begin{equation*}
 T^{\mu\nu}(p_2,p_1)T_{\mu\nu}(k_-,k_+)
 =32[(p_2\cdot k_-)(p_1\cdot k_+)+(p_2\cdot k_+)(p_1\cdot k_-)]
 =8(\hat u^2+\hat t^2).
\end{equation*}
再乘初态自旋平均 $1/4$、颜色平均 $1/N_c$ 以及 $g_{\rm em}^4Q_f^2/\hat s^2$，便得振幅平方。用 $z_\theta=\cos\theta$，
\begin{align*}
 \hat t&=-\frac{\hat s}{2}(1-z_\theta),\qquad
 \hat u=-\frac{\hat s}{2}(1+z_\theta),\\
 \frac{\dd\hat\sigma}{\dd z_\theta}
 &=\frac{\hbar^2}{32\pi\hat s}\overline{|\mathcal M|^2}
 =\frac{\pi\alpha^2\hbar^2Q_f^2}{2N_c\hat s}(1+z_\theta^2).
\end{align*}
由于 $\int_{-1}^{1}\dd z_\theta(1+z_\theta^2)=8/3$，角积分给\eqnrefs{eq:dy-born-partonic}。

将 Born 截面乘上两条运动学分布
\begin{equation*}
 \delta(\mathcal Q^2-sx_1x_2)
 \delta\!\left(Y-\frac12\ln\frac{x_1}{x_2}\right).
\end{equation*}
积分时需要的行列式为
\begin{equation*}
 \det\frac{\partial(sx_1x_2,\frac12\ln(x_1/x_2))}{\partial(x_1,x_2)}
 =\det\begin{pmatrix}sx_2&sx_1\\1/(2x_1)&-1/(2x_2)\end{pmatrix}=-s.
\end{equation*}
因此两次 delta 积分产生 $1/s$，而不是 $1/(x_1x_2s)$。部分子截面中的 $1/\hat s$ 则在支撑点变为 $1/\mathcal Q^2$，二者合并给\eqnrefs{eq:dy-born-hadronic}。其量纲为面积除以动量平方，与左边完全一致。
\end{derivation}

\begin{exbox}[海夸克不对称的近似读取]
在质子束打氢、氘靶的前向区域，若束流 $x_1$ 较大且由 $u$ 价夸克主导，靶端 $x_2$ 较小且由海夸克提供反夸克，则
\begin{equation*}
 \sigma_{pp}\propto\frac49u_p(x_1)\bar u_p(x_2),\qquad
 \sigma_{pn}\propto\frac49u_p(x_1)\bar d_p(x_2).
\end{equation*}
第二式使用同位旋近似 $\bar u_n=\bar d_p$。再把氘核近似成互不作用的质子加中子，得到
\begin{equation}
 \frac{\sigma_{pd}}{2\sigma_{pp}}
 \simeq\frac12\left[1+\frac{\bar d_p(x_2)}{\bar u_p(x_2)}\right].
 \label{eq:dy-sea-asymmetry-ratio}
\end{equation}
若同一分箱的比值为 $1.10$，这个近似对应 $\bar d_p/\bar u_p\simeq1.20$。它解释了轻子对为何能探测海夸克味结构。忽略的束流 $d$ 贡献由电荷权重 $1/4$ 及 $d_p/u_p$ 的乘积控制；束流反夸克、氘核束缚与卷积效应、高阶 QCD 和同位旋破缺也会改变提取结果。因此精确分析应回到完整味和及核 PDF，而不能把这个读数当作无模型测量。
\end{exbox}
